% !TeX program = latexmk
\documentclass[11pt,a4paper]{article}
\usepackage[utf8]{inputenc}
\usepackage[T1]{fontenc}
\usepackage[turkish]{babel}
\usepackage{lmodern}
\usepackage{geometry}
\geometry{margin=1in}
\usepackage{xcolor}
\usepackage{hyperref}
\hypersetup{
  colorlinks=true,
  linkcolor=blue,
  citecolor=teal,
  urlcolor=magenta,
  pdftitle={MCP Güvenliği},
  pdfauthor={Yusuf Talha Arabacı ve ark.}
}
\usepackage{graphicx}
\usepackage{tikz}
\usetikzlibrary{arrows.meta,positioning,fit}
\usepackage{csquotes}
\usepackage{booktabs}
\usepackage{enumitem}
\usepackage{microtype}
\usepackage[backend=biber,style=numeric,sorting=none]{biblatex}
\addbibresource{bibliography/references.bib}

\title{Model Bağlam Protokolü (MCP) Güvenliği:\\ Tehditler, Savunmalar ve Uygulama Rehberi}
\author{Yusuf Talha Arabacı ve ark.}
\date{Taslak: \today}

\begin{document}
\maketitle
\begin{abstract}
Bu çalışma, Model Bağlam Protokolü (MCP) güvenlik görünümünü; tehdit modelleri, saldırı sınıfları ve güvenli devreye alma için pratik savunmalarla birlikte ele almaktadır.
\end{abstract}

\section{Giriş}
% Intro\nThis paper investigates MCP security, contributions: (i) structured threat model, (ii) taxonomy of attacks, (iii) actionable defenses and deployment guidance.

\section{Arkaplan}
MCP abstracts tool and data access for LLM-enabled applications into a common protocol with defined roles: an MCP client (typically the LLM application), one or more MCP servers (brokers of tools and resources), and external tools/services. The protocol is often carried over JSON-RPC-like transports.

Key elements
- Client: issues requests, manages session context, and enforces UI/UX policies.
- Server: exposes tools and data sources through a unified interface and mediates execution.
- Tools/Services: concrete actions (e.g., filesystem, HTTP APIs, databases) with diverse risk profiles.

Trust boundary: clients should not implicitly trust servers or tools; servers should constrain tools; tools must be sandboxed or limited by design.

\begin{figure}[h]
\centering
\begin{tikzpicture}[
  comp/.style={rectangle,draw,rounded corners,align=center,minimum width=3.2cm,minimum height=1cm},
  arr/.style={-Latex}
]
\node[comp, fill=blue!5] (client) {MCP Client\\(LLM App)};
\node[comp, fill=green!5, right=3.2cm of client] (server) {MCP Server};
\node[comp, fill=orange!10, above right=0.8cm and 3.2cm of server] (tool1) {Tool A};
\node[comp, fill=orange!10, below right=0.8cm and 3.2cm of server] (tool2) {Tool B};
\draw[arr] (client) -- node[above]{Transport} (server);
\draw[arr] (server) -- (tool1);
\draw[arr] (server) -- (tool2);
\end{tikzpicture}
\caption{Simplified MCP topology and trust boundary.}
\label{fig:mcp-arch}
\end{figure}

Operational background
- Deployments vary from local developer environments (high privileges, low isolation) to cloud multi-tenant setups (stricter isolation and auditing).
- Misconfigurations (open servers, weak authn/z) significantly increase exposure \cite{MDPIElectronics3267}.


\section{Tehdit Modeli}
Varlıklar
- Gizli veriler: API anahtarları, kimlik bilgileri, müşteri/kurumsal veriler.
- Eylem bütünlüğü: araç çağrıları, dosya değişiklikleri, yürütme çıktıları.
- Kullanılabilirlik: istemci/sunucu çalışma süresi, oran ve bütçe sınırları.

Saldırganlar
- Ağ üzerinden dış aktör.
- Kötü niyetli/ele geçirilmiş araç veya sunucu sağlayıcı (tedarik zinciri).
- İç tehdit veya ele geçirilmiş işletim/uç nokta.

Varsayımlar
- Taşıma katmanı gizlilik/bütünlük sağlar (TLS), ancak uçlar güvenilir olmayabilir.
- Modeller halüsinasyon üretebilir; içerik veya araç meta verisi yoluyla prompt enjeksiyonu mümkündür.
- Araçların etki alanı farklıdır (dosya sistemi > salt-okunur HTTP gibi).

Güvenlik hedefleri
- Gizlilik: sırların araç yanıtları ve günlükler ile sızmaması.
- Bütünlük: yetkisiz/istenmeyen yan etkilerin önlenmesi.
- Kullanılabilirlik: kaynak kullanımını sınırla; saldırı altında güvenli bozulma.

Güven Sınırları
- İstemci $\leftrightarrow$ Sunucu: kimlik doğrulama, yetkilendirme, yetenek/kapasite mutabakatı.
- Sunucu $\leftrightarrow$ Araçlar: yalıtım ve politika uygulama.
- Kullanıcı $\leftrightarrow$ Model: prompt hijyeni, içerik filtreleme, kullanıcı onayı.


\section{Zafiyetler ve Saldırılar}
% Vulnerabilities and attacks: prompt/tool poisoning, plan injection, RCE, SSRF/DNS rebinding, authn/z gaps, supply chain.

\section{Savunmalar ve En İyi Uygulamalar}
% Defenses: isolation/sandbox, least privilege, validation, policy, observability, rate limiting; protocol-level proposals.

\section{İlgili Çalışmalar}
% Related work: agent frameworks, plugin ecosystems, API security.

\section{Sonuç}
We presented a concise threat model and an attack taxonomy for MCP deployments, along with pragmatic defenses and operational guidance. The central theme is disciplined trust management across client–server–tool boundaries, combined with isolation, least privilege, validation, and observability. Future work includes formal capability descriptors, server attestation mechanisms, and comparative evaluations of isolation strategies (e.g., containers vs. microVMs) under representative workloads.



\printbibliography
\end{document}
